\begin{enumerate}[label=\textbf{\arabic*)}, leftmargin=*]

    \item 
    O efeito Seebeck ocorre quando dois metais diferentes são unidos, formando uma junção, e a diferença de temperatura entre as junções gera uma tensão elétrica.
    O efeito Peltier ocorre quando uma corrente elétrica atravessa a junção de dois materiais diferentes, fazendo com que uma junção absorva calor e a outra libere calor, dependendo da direção da corrente.

    \item É um sensor de temperatura baseado no efeito Seebeck, ou seja, de acordo com a diferença de temperatura, gera tensão elétrica.

    \item 
        \begin{itemize}
            \item \textbf{Cobre-constantan (Tipo T):} São indicados para medições em ambientes excessivamente úmidos. Têm como principal desvantagem a baixa resistência à corrosão do cobre em temperaturas acima de 315\,°C;

            \item \textbf{Ferro-constantan (Tipo J):} São de baixo custo e versáteis, indicados para medições em ambientes não oxidantes devido ao ferro na composição da liga;

            \item \textbf{Cromel-alumen (Tipo K):} Apresentam boa resistência mecânica a altas temperaturas e são amplamente utilizados na indústria;

            \item \textbf{Cromel-constantan (Tipo E):} São adequados para medições em ambientes corrosivos.
        \end{itemize}

    \item O efeito piezoelétrico manifesta-se pelo aparecimento de uma diferença de potencial (tensão) entre faces opostas de um cristal quando este é deformado.
    No efeito piezoelétrico inverso, a aplicação de uma tensão elétrica sobre as duas faces de um cristal provoca a sua deformação mecânica.

    \item A temperatura de Curie representa a temperatura em que o sensor piezoelétrico perde suas propriedades piezoelétricas, ou seja, deixa de gerar tensão elétrica quando deformado.

    \item Um sensor piroelétrico detecta a variação de temperatura a partir da radiação infravermelha, gerando um sinal elétrico em resposta à mudança térmica.

    \item O efeito Hall é a geração de uma tensão elétrica perpendicular ao fluxo de corrente elétrica e a um campo magnético aplicado sobre um material condutor ou semicondutor.

    \item O sensor Hall unipolar só opera com um campo magnético de uma única polaridade (geralmente polo sul), enquanto o sensor Hall bipolar exige a alternância entre os campos sul e norte.

    \item O efeito magnetoresistivo é a variação da resistência elétrica do material de acordo com a intensidade e a direção do campo magnético aplicado.

    \item É a mudança de forma ou dimensões de um material ferromagnético quando ele é submetido a um campo magnético.

\end{enumerate}
